\documentclass{article}
\usepackage{amsmath}
\usepackage{amssymb}
\begin{document}
\begin{gather*}
\text {1. Prove. } 
\\ 
\exists x\in \mathbb{R}, x^3-2x+x=0
\\
\\
\text{Assume } x=0
\\
\text{Evaulate } x \text{ for } 0 \text{ in } x^3-2x+x=0
\\
0^3-2(0)+0=0
\\ 
0=0
\\ \\

\blacksquare \text{ } \exists x\in \mathbb{R}, x^3-2x+x=0 \ (\text{witness: }x=0)

\\ \\ \\ 

\text{2. Disprove. }
\\
\forall x\in \mathbb{Z}, \frac{x}{2}\in \mathbb{z}
\\ \\
\text{Assume }x=5
\\
\frac{5}{2} = 2.5
\\
2.5\notin \mathbb{Z}
\\ \\
\blacksquare \ \exists x \in \mathbb{Z}, \frac{x}{2} \notin \mathbb{Z} \ \text{(witness: } x=5)
\\ \\ \\ 

\text{3. Prove These Directly from the new definitions provided.}
\\ \\
\text{Definition A: } f(x) \text{ is increasing on domain } \mathbb{R} \text{ if and only if } \\ \forall x \in \mathbb{R}, \forall k > 0, f(x) < f(x+k)
\\ \\ 
\text{Show that the sum of increasing functions is increasing}

\\ \\ \text{Assume } f(x) \text{ is increasing.}
\\ \text{Applying the definition of an increasing function:} \\
\forall x \in \mathbb{R}, \forall k > 0, f(x) < f(x+k)
\\ \text{Assume the same for the function } g(x) \text{ and apply the same definition:} \\
\forall x \in \mathbb{R}, \forall k > 0, g(x) < g(x+k) \\ \\
\text{By the nature of addition} \\
f(x)+g(x) < f(x+k)+g(x+k) \ \forall k > 0
\\ \\
\blacksquare \ \forall x \in \mathbb{R}, \forall k > 0, f(x)+g(x) < f(x+k)+g(x+k)

\\ \\ \\ 

\text{Definition B: }f(x) \text{ is p-periodic on } \mathbb{R} \text{ if and only if} \\
\forall x \in \mathbb{R}, f(x+p) = f(x)
\\ \\ 
\text{Assume } f(x) \text{ is perdiodic, } \forall x \in \mathbb{R}, f(x+p) = f(x) \\
\text{Assume } g(x) \text{ is perdiodic, } \forall x \in \mathbb{R}, g(x+p) = g(x) \\
\text{Let }h(x) = f(x)+g(x) \\
\text{Then } h(x+p)=f(x+p)+g(x+p)
\\ \\
\text{Therefore }h(x+p)=h(x) \\
\text{Or} \\
f(x+p)+g(x+p) = f(x) + g(x) \\
\blacksquare \ \text{ which is p-peridiodic by the established definition.} \\ \\ \\

\text{Definition C: } f(x) \text{ is bounded above if and only if} \\
\exists K \in \mathbb{R}, \forall x \in \mathbb{R}, f(x) \le K \\ \\
\text{Prove: if } 𝑓(𝑥) \text{ is bounded above then } 2 𝑓(𝑥) \text{ is bounded above} \\ \\
\text{Let } 2K = L \text{ and }2f(x) = g(x)\\
\text{It follows that} \ 
\exists L \in \mathbb{R}, \forall x \in \mathbb{R}, g(x) \le L \\
\blacksquare \ \text{Therefore, } 2f(x) \ [g(x)] \text{ is bounded above by } 2K \ (L)
\end{gather*}
\end{document}
